\documentclass[11pt]{amsart}
\usepackage[OT1]{fontenc}
\renewcommand{\rmdefault}{cmr}
\usepackage[margin=1in]{geometry}
\usepackage{amsmath,amssymb,amsthm,booktabs,microtype}
\usepackage[hidelinks]{hyperref}
\newtheorem{theorem}{Theorem}
\newtheorem{lemma}[theorem]{Lemma}
\newtheorem{proposition}[theorem]{Proposition}
\newcommand{\Z}{\mathbb Z}
\newcommand{\Q}{\mathbb Q}
\newcommand{\rank}{\operatorname{rank}}
\title{The fourth integral Betti number of a torus-grid Rips complex}
\author{Kaiwen Jin}
\email{yui.ui.siki@gmail.com}
\date{October 14, 2026}
\subjclass[2020]{Primary 55N31; Secondary 05E45, 55U10}
\keywords{Vietoris--Rips complex, torus grid, integral homology, exact certificate}
\hypersetup{pdftitle={The fourth integral Betti number of a torus-grid Rips complex},pdfauthor={Kaiwen Jin}}

\begin{document}
\begin{abstract}
We compute the fourth integral Betti number of the closed radius-five
Vietoris--Rips complex on the thirteen-by-thirteen torus grid with its
periodic Manhattan metric. It is two, giving a counterexample to the
integer-rank formula in AIM Problem 531, which predicts twenty-six for this
grid. The proof separates a subcomplex with the homotopy type of the torus
from the simplices that cannot be lifted to the plane. A signed integer
chain certificate computes the resulting relative homology. Its nonzero
even boundary coefficients explain the difference from the published
mod-two data.
\end{abstract}
\maketitle

\section{The coefficient distinction}
For a periodic grid, the Vietoris--Rips construction can detect topology
beyond that of the underlying torus. Adams, Adetowubo, Barriga-Acosta, Feng
and Sterling~\cite{torus} determine several infinite families of torus-grid
Rips complexes and tabulate homology in the intermediate-scale range.
Their Question~8.4 concerns grids of side length $3k-2$ at scale $k$.
The tables use coefficients in $\Z/2$, whereas their section~3.6 defines
Betti numbers as ranks of integral homology. This distinction matters in
the integral formulation of AIM Problem~531~\cite{aim}.

For $n\geq3$, let $T_n=(\Z/n\Z)^2$, with
\[
 d_n((a,b),(c,e))=\delta_n(a,c)+\delta_n(b,e),\qquad
 \delta_n(a,c)=\min\{|a-c|,n-|a-c|\},
\]
where coordinates in the last formula lie in $\{0,\ldots,n-1\}$.
Let $K_n(r)$ consist of all subsets of $T_n$ of diameter at most $r$.
Thus equality with the scale is included. Write
$\beta_i(K)=\rank H_i(K;\Z)$.

The integral assertion in Problem~531 includes, for every integer $k\geq5$,
\begin{equation}\label{eq:prediction}
 \beta_4(K_{3k-2}(k))=6k-3+(-1)^k.
\end{equation}
The following single admissible case disproves that universal assertion.

\begin{theorem}\label{thm:main}
For the closed complex $K=K_{13}(5)$ with the metric above,
\[
 \rank H_4(K;\Z)=2.
\]
In particular, \eqref{eq:prediction} fails at $k=5$, where its right-hand
side is $26$.
\end{theorem}

The proof is computer-assisted. It uses an exact finite certificate, with
complete enumeration and integer elementary operations, together with a
geometric identification of a torus subcomplex. The certificate and a
separate reconstruction and replay program accompany the paper. The result
concerns the integral assertion at this one parameter; the mod-two pattern
in Question~8.4 remains a separate question.

\section{A subcomplex with the homotopy type of the torus}
From now on $n=13$ and the scale is five. Let
$\pi:\mathbb R^2\to\mathbb R^2/(13\Z)^2$ be the quotient map, and give
the plane and torus their Manhattan metrics. Define $M\subseteq K$ to
contain those simplices whose vertices have lifts to $\Z^2$ of diameter
at most five. Taking subsets preserves this property, so $M$ is a
subcomplex.

\begin{lemma}\label{lem:core}
There is a homotopy equivalence $M\simeq\mathbb R^2/(13\Z)^2$.
Consequently,
\[
 H_3(M;\Z)=H_4(M;\Z)=0,
 \qquad H_4(K;\Z)\cong H_4(K,M;\Z).
\]
\end{lemma}

\begin{proof}
Use the open radius-three torus balls centered at the 169 grid vertices.
They cover the torus, since every point has a grid vertex at distance at
most one. We first identify their nerve with $M$.

In the plane, set $u=x+y$ and $v=x-y$. For two points, Manhattan distance
equals $\max\{|\Delta u|,|\Delta v|\}$. A set of integer points has
diameter at most five precisely when both coordinate ranges have width
at most five. In that case the intervals
\[
 (\max u-3,\min u+3),\qquad
 (\max v-3,\min v+3)
\]
are nonempty, so its open radius-three balls have a common point.
Conversely, a common point gives pairwise distance less than six, which
is at most five for integer points.

Now take a nonempty intersection of torus balls, choose one grid center
and fix its plane lift $z_0$. Every point of the intersection has a unique
lift $w$ with $\|w-z_0\|_1<3$. For each other center choose a lift $z_i$
with $\|z_i-w\|_1<3$. Then $\|z_i-z_0\|_1<6$.
This lift of the center is independent of $w$: two different lifts are
at least thirteen apart, while two such choices would be less than twelve
apart. Thus the torus intersection is exactly the projection of the
fixed convex intersection $\bigcap_i B(z_i,3)$.
Projection is injective there, because it lies in $B(z_0,3)$.
Every nonempty intersection is therefore contractible, and the preceding
plane observation identifies the nerve with $M$.
The open-cover nerve theorem~\cite[Theorem~3.1]{torus} gives the claimed
homotopy equivalence.

The torus has a two-dimensional CW structure. Homotopy invariance, cellular
homology and the comparison with simplicial homology give the two
vanishing statements~\cite[Corollary~2.11 and Theorems~2.27,
2.35]{hatcher}. For completeness, the fourth-homology isomorphism can be
seen directly on quotient chains. A relative four-cycle lifts to a
four-chain whose boundary is a three-cycle in $M$; subtracting a four-chain
in $M$ with that boundary gives an absolute cycle. If an absolute cycle is
a relative boundary, subtract that boundary to leave a four-cycle in $M$,
which bounds in $M$. These operations prove surjectivity and injectivity.
\end{proof}

We also need an explicit way to recognize $M$ in a finite enumeration.
It is generated by the 338 projected closed diamonds of radius $5/2$
with centers
\begin{equation}\label{eq:diamonds}
 (a+\tfrac12,b),\qquad (a,b+\tfrac12),\qquad a,b\in\Z/13\Z.
\end{equation}
Indeed, for a lifted simplex choose the center in $(u,v)$ coordinates to
be $(\min u+5/2,\min v+5/2)$. Both coordinate ranges lie within distance
$5/2$ of this center. The corresponding $(x,y)$ center has one integral
and one half-integral coordinate, so it is of the form
\eqref{eq:diamonds}. Conversely, every such diamond has diameter five.
This is the scale-five instance of the planar facet geometry in
\cite[Proposition~4.2]{torus}; the argument above proves the membership
test used here.

\section{The signed relative chain calculation}
We use the relative simplicial chain groups
$C_q(K,M;\Z)=C_q(K;\Z)/C_q(M;\Z)$. Their basis consists of the simplices
of $K$ that are not in $M$. Label a vertex $(a,b)$ by $13a+b$, orient every
simplex by increasing labels, and order the basis lexicographically in
each degree. Its boundary is the alternating sum of faces, omitting those
in $M$.

\subsection{Complete enumeration}
Every edge of $K$ has a unique shortest plane displacement, because its
length is at most five and $13>10$. Fix a vertex of a clique as anchor and
lift its other vertices by these displacements. The clique lies in $M$
if and only if all pairwise distances of these lifts are at most five.
If the test fails, a triangle containing the anchor has nonzero total
edge displacement in $(13\Z)^2$. To see this, suppose every anchored
triangle has zero displacement sum. Then every difference between two
anchored lifts is the unique displacement of their edge and has length
at most five, contrary to failure of the test.

This gives a finite exhaustive enumeration: list all nonliftable triangles
through $(0,0)$, extend each by all cliques in its common-neighbor set,
translate by every element of $T_{13}$, and remove duplicates. Every
nonliftable simplex contains one of these translated triangles, and every
extension remains nonliftable. There is no degree cutoff in this
enumeration. It returns 24 origin triangles, at most eight possible
extension vertices per triangle, and no nonliftable simplex with more
than seven vertices.

An alternative reconstruction starts with the full neighbor graph and
recursively lists cliques through $(0,0)$ in increasing label order.
It tests membership in the diamonds~\eqref{eq:diamonds}, without using
the triangle criterion. A branch is skipped only when the union of the
current clique and all its possible extension vertices lies in one
diamond. Every descendant then belongs to $M$. Translating the retained
cliques gives exactly the same relative bases.

\medskip
\noindent\textbf{Finite verification procedure.}
The input is the closed metric graph on $T_{13}$ and the fixed cancellation
certificate. The output is its relative fourth-homology free rank.
\begin{enumerate}
\item Enumerate all noncore cliques by either reconstruction above.
\item Form the signed relative boundary matrices in the prescribed bases,
      and check that every consecutive boundary composition is zero.
\item Replay the recorded unit cancellations, checking each pivot and both
      adjacent chain identities.
\item Compare the complete residual integer matrix and count its nonzero
      columns; compute its rank modulo three.
\item Use the matching minor and column bounds to determine its rational
      rank, and subtract that rank from the residual fourth chain rank.
\end{enumerate}
The triangle construction considers at most $24\cdot2^8$ anchored
extension subsets and 169 translates per subset before deduplication.
The stored relative basis has 12168 simplices, each with at most seven
vertices, and each initial boundary has at most seven terms.
The recorded elimination has 6056 unit pivots. These are bounds and counts
for the fixed input, not a complexity assertion for arbitrary torus grids.

\subsection{Integral cancellation}
The following elementary fact accounts for every elimination step.

\begin{lemma}\label{lem:cancellation}
In a chain complex of finite free abelian groups, a boundary entry equal
to $1$ or $-1$ can be cancelled. If, after permuting bases, that boundary is
\[
 A=\begin{pmatrix}a&\beta\\ \gamma&D\end{pmatrix},\qquad a\in\{1,-1\},
\]
the remaining boundary is $D-\gamma a^{-1}\beta$. The resulting chain
complex has the same integral homology.
\end{lemma}

\begin{proof}
The integer invertible matrices
\[
 L=\begin{pmatrix}a^{-1}&0\\-\gamma a^{-1}&I\end{pmatrix},\qquad
 S=\begin{pmatrix}1&-a^{-1}\beta\\0&I\end{pmatrix}
\]
give $LAS=\operatorname{diag}(1,D-\gamma a^{-1}\beta)$.
Change the adjacent differentials by the corresponding inverse basis
changes. The chain identities force the coefficient in the distinguished
upper-boundary row and the distinguished lower-boundary column to vanish.
The changed complex therefore splits off a two-term summand
$\Z\xrightarrow{1}\Z$, whose homology is zero.
Deleting this summand proves the assertion.
\end{proof}

In a recorded step, let $r$ be the pivot row and $c$ the pivot column,
with entry $a=\pm1$. A column with entry $b$ in row $r$ is replaced by
itself minus $(b/a)$ times column $c$. Delete $r$ and $c$, together with
their corresponding column and row in the adjacent differentials.
Indices remain as labels for the changed bases. The replay program checks
that the adjacent directions vanish after the basis changes, as in the
lemma. In particular, a coefficient of absolute value two is never used
as a unit pivot.

\begin{proposition}\label{prop:finite}
The fixed certificate reduces $C_*(K,M;\Z)$ to a complex with nonzero
chain groups only in degrees $2,3,4$, of ranks $2,28,26$, respectively.
Its only nonzero differential is an integer matrix
$A_4:\Z^{26}\to\Z^{28}$ with exactly 24 nonzero columns and
\[
 \rank_{\Q} A_4=24.
\]
\end{proposition}

\begin{proof}
Complete reconstruction gives the initial chain ranks in
Table~\ref{tab:chains}. Signed boundary squares vanish over $\Z$.
The two enumerations agree on every basis simplex, and the independent
replay checks all 6056 pivots of the certificate using
Lemma~\ref{lem:cancellation}. The counts in the table leave exactly the
stated residual chain groups.

\begin{table}[ht]
\centering
\caption{Exact relative chain ranks and recorded unit cancellations.
A degree-$q$ cancellation removes one generator in degrees $q$ and $q-1$.}
\label{tab:chains}
\begin{tabular}{rrrr}
\toprule
Degree $q$ & Initial rank & Degree-$q$ cancellations & Residual rank\\
\midrule
2 & 1352 & 0 & 2\\
3 & 4394 & 1350 & 28\\
4 & 4394 & 3016 & 26\\
5 & 1690 & 1352 & 0\\
6 & 338 & 338 & 0\\
\bottomrule
\end{tabular}
\end{table}

The certificate contains the entire residual $28\times26$ matrix.
All its nonzero entries are $2$ or $-2$, and exactly 24 columns are
nonzero. This bounds its rational rank above by 24. Exact row reduction
modulo three gives rank 24; consequently some $24\times24$ integral
minor is nonzero modulo three, hence nonzero over $\Q$. This gives the
matching lower bound and proves the proposition.
\end{proof}

\begin{proof}[Proof of Theorem~\ref{thm:main}]
After the homology-preserving cancellations, the fifth chain group is
zero, so $H_4(K,M;\Z)$ is the integral kernel of $A_4$.
The rational span of this kernel equals $\ker(A_4:\Q^{26}\to\Q^{28})$:
clear denominators in any rational kernel vector. Thus its free rank is
$26-24=2$. Lemma~\ref{lem:core} transfers this rank to $H_4(K;\Z)$.
Since $13=3\cdot5-2$ and $6\cdot5-3+(-1)^5=26$, this is an admissible
counterexample to \eqref{eq:prediction}.
\end{proof}

\section*{Reproducibility}
The accompanying files contain a deterministic generator,
\path{src/torusgrid_certificate.py}, the independent reconstruction and
replay program \path{src/verify_certificate.py}, and the proof-critical
input \path{numerical_case/certificates/k5_certificate.json}.
Python 3.10 or later with its standard library suffices. From the artifact
root, run
\begin{verbatim}
python3 src/verify_certificate.py \
  numerical_case/certificates/k5_certificate.json
python3 -m unittest discover -s test -v
python3 numerical_case/generate_certificate.py \
  --output /tmp/k5-regenerated.json
\end{verbatim}
The generator requires a fresh output path. Its output contains no
timestamps or platform-dependent values and can be compared directly with
the fixed certificate. The verifier imports no generator code. It checks
the finite calculation; the geometric core and the transfer to the original
complex are proved above. The tests include a four-sphere boundary and a
triangulation of the real projective plane, checking in particular that even
boundary coefficients are preserved. The zero reduction of $A_4$ modulo
two accounts for the larger fourth dimension in the mod-two calculation.

\section*{Automated assistance}
OpenAI Codex was used for proof exploration, exact verification software,
source checking and manuscript preparation. The argument and certificate
are submitted for independent mathematical review.

\begin{thebibliography}{9}
\bibitem{torus}
H.~Adams, A.~Y.~Adetowubo, H.~Barriga-Acosta, Z.~Feng and J.~Sterling,
\emph{Vietoris--Rips Complexes of Torus Grids},
Mediterranean Journal of Mathematics \textbf{22} (2025), article 173,
\href{https://doi.org/10.1007/s00009-025-02945-9}{doi:10.1007/s00009-025-02945-9}.
Author version: \href{https://arxiv.org/abs/2502.07134v2}{arXiv:2502.07134v2}.
\bibitem{hatcher}
A.~Hatcher, \emph{Algebraic Topology}, Cambridge University Press, 2002.
\bibitem{aim}
\emph{AIM problem collection}, GitHub repository \texttt{MColbrook/AIM}, Problem 531,
``Homology at the second transition of torus-grid Rips complexes'',
revision \texttt{a153f469ee0ea64e876b2ad97c79d252c19a3e22},
\href{https://github.com/MColbrook/AIM/blob/a153f469ee0ea64e876b2ad97c79d252c19a3e22/problems/531-torus-grid-transition-homology.md}{immutable problem statement}.
\end{thebibliography}
\end{document}
